\documentclass[11pt]{article}

\usepackage[margin=1in]{geometry}
\usepackage{amsmath,amssymb}
\usepackage{graphicx}
\usepackage{booktabs}
\usepackage{siunitx}
\usepackage[hidelinks]{hyperref}
\usepackage{xcolor}
\usepackage{caption}

\graphicspath{{figures/}}

\newcommand{\erf}{\ensuremath{\operatorname{erf}}}
\newcommand{\ReLU}{\ensuremath{\operatorname{ReLU}}}
\newcommand{\Sab}{\ensuremath{S_{\mathrm{ab}}}}
\newcommand{\Sinc}{\ensuremath{S_{\mathrm{inc}}}}
\newcommand{\Teff}{\ensuremath{T_{\mathrm{eff}}}}
\newcommand{\Tavg}{\ensuremath{\bar{T}}}
\newcommand{\khat}{\ensuremath{\hat{k}}}
\newcommand{\nhat}{\ensuremath{\hat{n}}}
\newcommand{\rr}{\ensuremath{\mathbf{r}}}
\newcommand{\Rbar}{\ensuremath{\bar{R}}}

\title{How exact can the surface dosimetry get?\\
\large Inter-body reflection and polarization, tested on Thelonious}
\author{AEGIS self-shadowing study}
\date{\today}

\begin{document}
\maketitle

\begin{abstract}
The companion note verified the diffraction-smoothed visibility gate. This
note asks how much further the surface-dosimetry law can be pushed toward
the exact electromagnetic answer, and tests every claim on the real
Thelonious phantom mesh (\num{23826} triangles). Two refinements are
candidates: inter-body re-illumination and polarization. We find, first,
that the monograph's omission of inter-body reflection is essentially
correct for body-averaged exposure (enhancement $C\approx 1.05$--$1.08$),
but its specular estimate is optimistic: a forward specular-bounce test
gives a recapture fraction \emph{comparable} to the diffuse bound, not the
assumed one-third of it, and the deepest concavity reaches a local
$C\approx 1.85$. Second, and more consequentially, polarization is the
largest unexploited source of inaccuracy in the incoherent pipeline:
switching the incident wave between horizontal and vertical polarization
changes the whole-body absorbed power by \SI{25}{\percent}, the local
effective transmission deviates from its unpolarized value by up to
\SI{92}{\percent}, and \SI{44}{\percent} of the lit body lies at grazing
incidence ($\theta>\SI{60}{\degree}$) where the TE/TM splitting peaks.
The diffraction gate inherits a polarization signature of its own: for the
smooth convex human body the hard (TM) creeping wave decays
$2.3\times$ slower than the soft (TE) one, a \SI{6.4}{\dB} polarization
split \SI{15}{\degree} into the geodesic shadow. We conclude that the
highest-value exactness upgrade is not more elaborate diffraction but
carrying the true polarization state through the incoherent levels, with a
soft/hard gate as the natural and inexpensive refinement.
\end{abstract}

\section{Scope}

The production law is the first-bounce, polarization-averaged statement
\begin{equation}\label{eq:base}
  \Sab(\rr) = \Sinc\,T_0\,\ReLU\!\bigl(\mu(\rr)\bigr)\,O(\rr,\khat),
  \qquad \mu = \nhat\cdot(-\khat),
\end{equation}
with the visibility $O$ developed and verified in the companion note. Three
physical effects sit outside~\eqref{eq:base}: (i) one body part reflecting
power onto another (inter-body reflection); (ii) the dependence of the
absorbed power on the incident polarization, which the incoherent levels
0--6 replace with an angle-averaged transmission; and (iii) the
polarization dependence of the diffraction that smooths the shadow edge.
This note quantifies all three on the Thelonious mesh and identifies where
the accuracy is actually being lost. All numbers below are reproduced by
\texttt{studies/self\_shadowing/exactness\_study.py}.

\section{Inter-body reflection: small, but not for the reason assumed}

\paragraph{Theory.} At an illuminated point a fraction
$\Rbar = 1-T_0 \approx 0.46$ of the incident power is reflected. If a
fraction $f(\rr)$ of that reflected power strikes the body again, summing
the geometric series of bounces enhances the absorbed power by
\begin{equation}\label{eq:C}
  C(f) = \frac{1}{1-\Rbar\,f}.
\end{equation}
Under diffuse (Lambertian) reflection the recaptured fraction is bounded by
the cosine-weighted hemisphere blocked by the body, $f(\rr)\le 1-\eta(\rr)$,
where $\eta$ is exactly the ambient-occlusion exposure fraction AEGIS
already computes. The monograph argues this is negligible through a
self-compensation mechanism: $C$ is largest precisely where $\eta$ is
smallest, so a deep concavity sees less direct power but recaptures more,
and the product $\eta\,C$ stays close to~$\eta$.

\paragraph{Test.} We compute $\eta$ on Thelonious (\num{128} cosine-weighted
rays per triangle, production BVH), form the diffuse recapture map
$f = 1-\eta$ and the enhancement~\eqref{eq:C}, and separately fire the
specular bounce ray $\hat r = \khat - 2(\khat\cdot\nhat)\nhat$ from every
lit triangle over a sweep of \num{24} incidence directions to measure the
true specular recapture fraction. Results are summarized in
Table~\ref{tab:interbody} and Fig.~\ref{fig:interbody}.

\begin{table}[!htb]
\centering
\caption{Inter-body reflection on Thelonious. The diffuse bound and the
body-averaged enhancement confirm the effect is small; the specular
recapture and the worst-case concavity show the monograph's specular
estimate is optimistic.}
\label{tab:interbody}
\begin{tabular}{lcc}
\toprule
Quantity & This study & Monograph \\
\midrule
Diffuse recapture $f_{\mathrm{global}}$ & 0.135 & $\approx 0.09$ \\
Body-averaged enhancement $C_{\mathrm{global}}$ & 1.077 & $\approx 1.04$ \\
$\eta$-weighted enhancement & 1.052 & $\lesssim 1.02$ \\
Specular recapture $f_{\mathrm{spec}}/f_{\mathrm{diffuse}}$ & 1.08 & $\approx 0.3$ \\
Worst-concavity diffuse $C$ & 1.85 & --- \\
\bottomrule
\end{tabular}
\end{table}

\begin{figure}[!htb]
\centering
\includegraphics[width=\textwidth]{inter_body}
\caption{Left: area-weighted distribution of the diffuse recapture fraction
$f=1-\eta$ on Thelonious; the body average is \num{0.135}. Right: the
self-compensation mechanism. Effective exposure $\eta\,C$ stays close to the
no-reflection line $\eta$ for both the diffuse bound (red) and the specular
estimate (blue); deviations grow only in deep concavities (small $\eta$).}
\label{fig:interbody}
\end{figure}

\paragraph{Verdict.} The self-compensation argument holds: the
body-averaged enhancement is \SIrange{5}{8}{\percent}, well inside the
diffraction and tissue-property error bars, so omitting inter-body
reflection from the body-averaged law is justified. Two honest caveats
emerge from the data. (1) The diffuse recapture fraction we measure
(\num{0.135}) is roughly \SI{50}{\percent} above the monograph's quoted
\num{0.09}; the difference is pose- and resolution-dependent but does not
change the conclusion. (2) The specular recapture is \emph{not}
one-third of the diffuse bound as assumed. In concavities the mirror
direction points back at the opposing surface, so the specular bounce is
recaptured about as often as the diffuse hemisphere is blocked
($f_{\mathrm{spec}}\approx 1.08\,f_{\mathrm{diffuse}}$). The specular
enhancement is therefore close to the diffuse $C\approx 1.05$--$1.08$
rather than the $C\approx 1.01$ the monograph states, and the deepest
concavity reaches $C\approx 1.85$. This matters for spatial maps in tight
gaps (between the legs, under the arms), not for compliance, which is set
by the most-illuminated surface. Recommendation: keep inter-body reflection
out of the default law, but correct the specular claim in the error budget
and expose it as an optional first-bounce pass for spatial-map fidelity.

\section{Polarization: the largest unexploited inaccuracy}

\paragraph{Theory.} A plane wave is fully polarized. At a surface point the
incidence plane is spanned by $\khat$ and $\nhat$, with local TE and TM unit
vectors $\hat e_s = (\khat\times\nhat)/|\khat\times\nhat|$ and
$\hat e_p = \hat e_s\times\khat$. The electric field decomposes as
$\mathbf{E}_0 = a_s\hat e_s + a_p\hat e_p$, and the exact effective
power-transmission coefficient is
\begin{equation}\label{eq:Teff}
  \Teff(\rr) = |e_s|^2 T_s(\theta) + |e_p|^2 T_p(\theta)
  = \Tavg(\theta) + \tfrac{1}{2}\,q(\rr)\,\Delta T(\theta),
\end{equation}
with $\Tavg = \tfrac12(T_s+T_p)$, $\Delta T = T_p - T_s$, and the local
TM excess $q = |e_p|^2 - |e_s|^2 \in[-1,1]$. The exact absorbed power density
is $\Sab = \Sinc\,\Teff\,\ReLU(\mu)$. The incoherent levels 0--6 set
$\Teff = \Tavg$ (level~3) or apply a single scalar $q$ knob (level~4), so
they cannot represent the spatial pattern of $q(\rr)$ that a real polarized
source imposes. The splitting $\Delta T$ is the lever: it vanishes at normal
incidence and grows steeply with angle.

\begin{figure}[!htb]
\centering
\includegraphics[width=0.46\textwidth]{deltaT_angle}
\caption{Polarization splitting $\Delta T(\theta) = T_p - T_s$ for skin at
three bands. The splitting vanishes at normal incidence and exceeds
\num{0.77} by \SI{75}{\degree}: at grazing, TM transmits almost completely
while TE is strongly reflected.}
\label{fig:deltaT}
\end{figure}

\paragraph{Test.} For a horizontally propagating plane wave ($\khat=+\hat x$)
on the upright Thelonious phantom we compute the per-triangle incidence
angle, $T_s,T_p,\Tavg$, the local TE/TM basis, and the integrated absorbed
power for vertical polarization, horizontal polarization, and the
unpolarized assumption. Results in Table~\ref{tab:pol},
Figs.~\ref{fig:deltaT}--\ref{fig:bodymap}.

\begin{table}[!htb]
\centering
\caption{Polarization on Thelonious, wave from $+\hat x$, skin at
\SI{28}{\GHz}. The incoherent pipeline uses the ``unpolarized'' column and
misses the rest.}
\label{tab:pol}
\begin{tabular}{lc}
\toprule
Quantity & Value \\
\midrule
$\Delta T$ at \SI{75}{\degree} (28\,GHz) & 0.77 \\
Whole-body absorbed power, vertical / unpolarized & 0.87 \\
Whole-body absorbed power, horizontal / unpolarized & 1.13 \\
Horizontal $-$ vertical (relative to unpolarized) & \SI{25}{\percent} \\
Max local $|\Teff-\Tavg|/\Tavg$ & \SI{92}{\percent} \\
Lit area at grazing ($\theta>\SI{60}{\degree}$) & \SI{44}{\percent} \\
\bottomrule
\end{tabular}
\end{table}

\begin{figure}[!htb]
\centering
\includegraphics[width=\textwidth]{polarization_body}
\caption{Left: area-weighted distribution of lit-surface incidence angle.
Almost half the illuminated body is at grazing ($\theta>\SI{60}{\degree}$),
where polarization dominates. Right: the resulting per-triangle deviation
of the effective transmission from its unpolarized value for vertical
polarization, reaching $\pm\SI{90}{\percent}$.}
\label{fig:polbody}
\end{figure}

\begin{figure}[!htb]
\centering
\includegraphics[width=0.7\textwidth]{polarization_bodymap}
\caption{Where polarization matters, on the body. Left: relative deviation
of $\Teff$ from the unpolarized value (vertical polarization). Near-vertical
surfaces (torso, limbs) face the horizontal wave at grazing and transmit
\emph{less} than the average for vertical polarization (blue); near-horizontal
surfaces (shoulders, head crown, insteps) transmit more (red). Right: the
incidence-angle map that drives it.}
\label{fig:bodymap}
\end{figure}

\paragraph{Verdict.} Polarization is the dominant unexploited inaccuracy in
the incoherent pipeline. A \SI{25}{\percent} swing in whole-body absorbed
power between the two linear polarizations, and a local effective-transmission
error approaching a factor of two, both exceed the diffraction and
tissue-property error bars that justified dropping inter-body reflection.
The effect is structural, not incidental: an upright body illuminated near
the horizontal presents \SI{44}{\percent} of its lit area at grazing, exactly
where $\Delta T$ is largest. The exact law~\eqref{eq:Teff} is already in the
monograph and is implemented in the coherent levels~7--8; the gap is that
levels 0--6 discard the polarization state. Carrying a single Stokes vector
(or the $q(\rr)$ field) through the incoherent kernels closes it at the cost
of two extra dot products per triangle and is, in our view, the single
highest-value exactness upgrade available.

\section{Polarization in the diffraction gate: soft and hard}

The same grazing region that maximizes $\Delta T$ is where the diffraction
gate operates, and diffraction is itself polarization-dependent. For a
perfectly conducting smooth convex body, the deep-shadow surface field is a
sum of creeping-wave modes $\propto \exp(i t_m\,\xi)$ in the Fock arc-length
parameter $\xi = m\,(s/R)$, with $m=(kR/2)^{1/3}$ and $s$ the geodesic depth
past the terminator on a feature of radius~$R$. The modal eigenvalues are
$t_m = q_m\,e^{i\pi/3}$, so each mode decays at rate
$\operatorname{Im} t_m = q_m\sin(\pi/3) = \tfrac{\sqrt3}{2}q_m$ per unit
$\xi$. The crucial point is that $q_m$ differs by polarization:
\begin{equation}
  \text{soft (TE, Dirichlet): } q_1 = 2.338\ (\text{zeros of } \mathrm{Ai}),
  \qquad
  \text{hard (TM, Neumann): } q_1 = 1.019\ (\text{zeros of } \mathrm{Ai}').
\end{equation}
The hard (TM) creeping wave therefore decays $2.295\times$ slower than the
soft (TE) one. On a \SI{5}{\cm} limb at \SI{28}{\GHz} ($m=2.45$), this is a
\SI{6.4}{\dB} polarization split only \SI{15}{\degree} of arc into the
geodesic shadow (Fig.~\ref{fig:softhard}). A scalar shadow gate cannot
represent this; a gate that carries the soft/hard label can, at no extra
baked storage, because the geometric inputs ($R$, $s$, $\xi$) are shared and
only the modal constant changes. (The human body is a lossy dielectric, not
a perfect conductor, so the exact constants are modified by the surface
impedance; the PEC values bracket the true split and preserve the sign and
$2.3\times$ asymmetry.)

\begin{figure}[!htb]
\centering
\includegraphics[width=\textwidth]{soft_hard}
\caption{Soft/hard creeping-wave split on a \SI{5}{\cm} limb at
\SI{28}{\GHz}. Left: surface field vs geodesic depth into shadow for the two
polarizations; the hard (TM) field penetrates much deeper. Right: the
polarization split grows past \SI{6}{\dB} within \SI{15}{\degree} of arc.}
\label{fig:softhard}
\end{figure}

This also resolves a sign question. The companion note found the scalar
erf gate \emph{over}-predicts deep-shadow leakage relative to a knife edge
(Gaussian vs $\nu^{-2}$ tail), but a knife edge is the wrong canonical
problem for a smooth body. The correct smooth-convex tail is the creeping
wave, which decays \emph{exponentially} in $\xi$, faster than either. So on
the body the scalar gate over-predicts leakage for both polarizations, and
most so for TE; the soft/hard refinement corrects the rate and splits it.

\section{Synthesis and recommendations}

\begin{enumerate}
\item \textbf{Carry polarization through levels 0--6 (highest value).} Use
the exact $\Teff = \Tavg + \tfrac12 q\,\Delta T$ with the true per-triangle
$q(\rr)$ from the source polarization, not a scalar knob. Cost: two dot
products per triangle. Payoff: removes a \SI{25}{\percent} whole-body and
up to \SI{92}{\percent} local error. This is already exact theory and is
the form the coherent levels use.
\item \textbf{Make the diffraction gate soft/hard.} Replace the single
scalar transition with a polarization-labelled creeping-wave gate using the
soft and hard Airy constants. The bake is unchanged (same geometric $R$,
$\xi$); only a constant switches. Payoff: a \SI{6.4}{\dB} polarization
signature in the shadow, and the physically correct exponential smooth-body
tail rather than the knife-edge surrogate.
\item \textbf{Keep inter-body reflection optional.} The body-averaged
enhancement ($\le\SI{8}{\percent}$) does not warrant inclusion in the default
law, but the error budget should record the specular recapture as
\emph{comparable to} the diffuse bound (not one-third of it) and the
worst-concavity $C\approx 1.85$. A single first-bounce pass, reusing the
visibility BVH, suffices for spatial-map fidelity in concavities if needed.
\end{enumerate}

\section{Verification suite}

Table~\ref{tab:tests} lists the checks that pin these results and that
should run in CI to keep them true. The numerical oracles are exact; the
physical checks are regression tests against the values reported here.

\begin{table}[!htb]
\centering
\caption{Verification tests for the exactness study.}
\label{tab:tests}
\begin{tabular}{lll}
\toprule
Test & Asserts & Oracle / tolerance \\
\midrule
$\Teff$ at normal incidence & $\Teff = T_0$ for any polarization & exact \\
$\Delta T(0)=0$ & no splitting at normal incidence & exact \\
$q\in[-1,1]$ & TM excess bounded & exact \\
$\Teff$ unpolarized average & $\langle\Teff\rangle_\phi = \Tavg$ & exact \\
Energy conservation & $T_s,T_p\in[0,1]$, $T+R=1$ lossless limit & exact \\
Fresnel vs monograph & $T_0(\text{skin},28\,\text{GHz})=0.539$ & $<1\%$ \\
Recapture bound & $f \le 1-\eta$ everywhere & exact \\
Self-compensation & $\eta\,C \le 1$ and monotone in $\eta$ & exact \\
Airy constants & $q_1^{\text{soft}}=2.338$, $q_1^{\text{hard}}=1.019$ & $<10^{-3}$ \\
Soft/hard ratio & $\alpha_{\text{soft}}/\alpha_{\text{hard}} = 2.295$ & $<10^{-3}$ \\
Thelonious regression & whole-body H/V ratio, grazing fraction & $\pm 2\%$ \\
\bottomrule
\end{tabular}
\end{table}

\section{Conclusion}

Pushing the surface law toward the exact electromagnetic answer, the biggest
win is not more elaborate diffraction or inter-body radiosity, both of which
the data show to be modest. It is polarization. The exact polarization-aware
transmission is already derived and already used in the coherent levels;
extending it to the incoherent levels removes the single largest error in
the pipeline at trivial cost, and a soft/hard creeping-wave gate carries the
same polarization distinction into the shadow edge where the body spends
\SI{44}{\percent} of its illuminated area. Inter-body reflection stays
optional, with its specular contribution corrected upward to match the
measured recapture.

\end{document}
